\usepackage{indentfirst}
\usepackage{float}
\floatplacement{figure}{H}

\usepackage{amsmath, amssymb}
\usepackage{mathtools}
\usepackage{bm}

% *** ОПРЕДЕЛЯЕМ МАТЕМАТИЧЕСКИЕ ШРИФТЫ ДО PLEX ***
\DeclareMathAlphabet{\mathrm}{OT1}{lmr}{m}{n}
\DeclareMathAlphabet{\mathbf}{OT1}{lmr}{bx}{n}
\DeclareMathAlphabet{\mathit}{OT1}{lmr}{m}{it}
\DeclareMathAlphabet{\mathcal}{OT1}{lmr}{m}{it}
\SetMathAlphabet{\mathrm}{bold}{OT1}{lmr}{bx}{n}
\SetMathAlphabet{\mathbf}{bold}{OT1}{lmr}{bx}{n}

\IfFileExists{plex-otf.sty}{
  % Теперь загружаем Plex - он увидит, что математика уже определена
  \usepackage[
    RM={Scale=0.94},SS={Scale=0.94},SScon={Scale=0.94},TT={Scale=MatchLowercase,FakeStretch=0.9},
    DefaultFeatures={Ligatures=Common}
  ]{plex-otf}
  
  % Дополнительная поддержка Unicode символов
  \usepackage{newunicodechar}
  \newunicodechar{𝑑}{d}
  \newunicodechar{𝑢}{u}
  \newunicodechar{𝛼}{\alpha}
  \newunicodechar{𝑡}{t}
  
  % Явная загрузка JuliaMono
  \setmonofont{JuliaMono}[
    Path = /home/glebq/.local/share/fonts/,
    Extension = .ttf,
    UprightFont = *-Regular,
    BoldFont = *-Bold,
    ItalicFont = *-RegularItalic,
    BoldItalicFont = *-BoldItalic
  ]
}{
  %% TinyTeX
  \usepackage{libertine}
  \usepackage{newunicodechar}
  \newunicodechar{𝑑}{d}
  \newunicodechar{𝑢}{u}
  \newunicodechar{𝛼}{\alpha}
  \newunicodechar{𝑡}{t}
}
